% ===== §18 =====
\section{The Description of Private-Cultural Phenomena}
\label{p24:sec:description}

\noindent
The description begins with the phenomena themselves, and they are so familiar that their strangeness is easily missed. Any relation of sufficient depth generates, over time, a texture of forms that belong to it and to no one else, and the first task is simply to see this texture for what it is: a culture, in the full sense the first part gave the word, generated between two people and bearing every mark the definition named, only at the smallest scale and with the greatest intimacy. The phenomena are small and are usually beneath notice precisely because they are so intimate, but catalogued plainly they compose an unmistakable cultural whole.

\subsection*{The inventory}

There is, first, a \emph{private lexicon}: words that mean, between these two, what they mean nowhere else. A word coined for a thing that has no public name; a public word bent to a private sense; a mispronunciation kept on purpose until it becomes the right pronunciation; a name for a mood, a look, a recurring situation, that the two of them share and no dictionary holds. There are \emph{private references}: the compressed allusion to a shared past event, a place, a person, a disaster, a joke, that carries in a word or a glance a whole history the two of them lived and no one else was present for. There are \emph{inside jokes}, which are private references raised to the pitch of the comic, a charge that discharges in laughter for exactly two people and is inert for any third. There are \emph{private names}, the appellations the two use for each other that no one else uses and that would sound wrong, even faintly transgressive, in another mouth---the diminutive, the epithet, the animal or the object that has become, between them, a name. There are \emph{micro-rites}: the small repeated ceremonies of a shared life, the particular way a departure or a return or a nightfall is marked, the gesture that must be made or the phrase that must be said, whose omission is felt as a wrongness though no external rule requires it. And there are \emph{understandings}: the wordless coordinations, the shared silences that are not empty but full of a determinate content both parties read, the glance across a room that communicates a complete judgement, the gesture that means, to these two, a paragraph. Taken together these are not a scattering of habits; they are the imaginary, symbolic, and real substance of a culture---the images and identifications, the coded signs, and the unsayable ``between''---in the concentrated form the dyad alone produces.

\subsection*{How they are given: the first- and second-person structure}

What is most distinctive about these phenomena is not their content but the \emph{structure of their givenness}, and it is here that the phenomenology begins to disclose the mechanisms the later sections will name. These forms are given only in the first- and second-person: they exist as, and only as, \emph{our} idiom, spoken by \emph{me} to \emph{you} and understood by you as from me. They have no third-person mode of existence in which they could be surveyed from outside without loss, because their meaning is not a content lodged in them but a charge that exists only in the circuit between the two who share them. This is not true of macro-cultural forms in the same way. A national custom, a public rite, a widely shared idiom has a third-person existence: an outsider can learn it, describe it, come to understand what it means to those who practise it, and though something of the lived sense may escape the outsider, the form itself is there to be studied in the third person, a public object. The private forms of an intimate culture are different in kind. There is no third-person standpoint from which the charge of an inside joke can be observed, because the charge \emph{is} the circuit between the two, and to stand outside the circuit is not to see the charge imperfectly but not to be in its presence at all. The outsider does not understand the private form less well than the two; the outsider is simply not one of the two for whom alone it is a form. This is why the phenomena are given in the first and second person irreducibly: they are addressed, they live in address, and an addressed meaning surveyed from the position of a non-addressee has ceased to be the thing it was.

\subsection*{The temporality of the forms}

One further feature of the phenomena must be marked before their knowing is taken up, because it will matter for the mechanism of sedimentation. These forms have a distinctive \emph{temporality}: each of them was, at some moment, \emph{new}, and became, over some stretch of shared time, \emph{settled}, and the passage from the new to the settled is short enough, in the dense medium of an intimate life, to be within living memory of both parties. A couple can often recall, with a precision impossible for any macro-cultural form, the occasion on which a private word was coined, the event a private reference compresses, the first time a micro-rite was performed before it was a rite. The forms carry their own origins with them; they are sedimented, but the sediment is recent enough that its laying-down is remembered. This is the phenomenological trace of a mechanism the later section on sedimentation will describe: that intimate culture emerges at a speed which brings the moment of emergence itself within the horizon of experience, so that the couple are not only the bearers of their culture but the witnesses of its coming-to-be. With the phenomena thus described---their inventory, their first- and second-person givenness, their remembered temporality---the question can be asked that the epistemology of this culture turns on: how, and by whom, such forms are known.
